% ============================================================
% Section 1: Introduction
% Negative diagnostic study: frozen graph-vs-MLP family selection
% ============================================================

\section{Introduction}
\label{sec:intro}

Graph Neural Networks (GNNs) can exploit relational information that a feature-only model cannot access, but using the graph is not uniformly beneficial~\cite{kipf2017gcn,velivckovic2018gat,hamilton2017graphsage}. This creates a practical decision before deployment: should a node-classification task use a graph-model portfolio or a feature-only multilayer perceptron (MLP) portfolio? Training and validating both portfolios can answer the question retrospectively, but a cheap diagnostic would be valuable when training and tuning the graph-model portfolio is costly. The graph is already available in this setting because both the candidate models and the structural diagnostics use it.

Existing work shows why this decision is difficult. Homophily alone does not characterize GNN performance, and recent metrics incorporate class-conditional, feature, structural, or neighborhood-distribution information~\cite{luan2023whendo,platonov2023characterizing,zheng2024trihom,pmlr-v235-wang24u}. Other work treats graph-model selection as a cross-dataset meta-learning problem~\cite{park2023metagl,park2023glemos}. These studies establish that graph characteristics can correlate with performance and that learned selectors can rank graph models. They do not, however, determine whether a fixed set of transparent statistics, with label-dependent graph quantities restricted to the training split, is reliable enough to support an explicit graph-versus-MLP family action under matched trivial and validation-based policies.

Correlation is not the same as a safe decision. A statistic can correlate with a performance gap yet still produce costly errors at a fixed threshold; a high selection accuracy can also be misleading when one action dominates the benchmark. We therefore evaluate diagnostics as policies with three outputs---graph, MLP, or abstain---and measure oracle-portfolio regret, coverage, and selection accuracy. The target selects graph only when the tuned graph portfolio exceeds the tuned MLP by more than a predeclared one-percentage-point margin; all smaller gaps map to MLP. Regret remains primary because graph is the target on 89 of the 110 evaluation units.

We conduct a frozen prospective evaluation on newly generated confirmatory records rather than redesign the rule after observing their comparative outcomes. Before generating those records, we fix the candidate datasets, ten seeds, 60/20/20 splits, seven architectures, four tuning trials per architecture, diagnostic label scope, thresholds, tie handling, MLP fallback for abstention, and dataset-level inference. Historical aggregate results had been inspected during earlier exploratory work, but they are excluded from the confirmatory benchmark. A feasibility-aborted v1 run reveals that Texas has a singleton class and cannot support the frozen per-class three-way split. Before any comparative benchmark was assembled or evaluated, an outcome-independent eligibility amendment excludes Texas, assigns a new run identifier, and restarts all remaining datasets from zero without changing the scientific settings. The final 11-dataset artifact contains 770 model records and 110 train-only diagnostic records with no failed evaluation units. The combined diagnostic is compared on identical units with always-graph, always-MLP, random 50/50, homophily-only, degree-only, homophily-plus-degree, two-hop-only, and direct validation-selection policies.

The prospective result is negative and operationally consequential. The combined diagnostic reaches 55.5\% selection accuracy at 68.2\% coverage and incurs 7.46 percentage points of full-set regret. Always-graph reaches 80.9\% accuracy with 0.26 points of regret, while validation selection reaches 91.8\% accuracy with 0.22 points of regret. None of the eight predeclared comparisons establishes an advantage for the combined rule after Holm correction. The evaluated diagnostics therefore provide no stable incremental value for graph-vs-MLP selection under this protocol.

This failure does not imply that graph structure is irrelevant. In a separate frozen intervention, degree-preserving randomization reduces GCN test accuracy by an average of 45.5 points on Cora, 29.9 on CiteSeer, and 17.3 on PubMed. The intervention preserves every node's degree but jointly changes class mixing, connectivity patterns, and higher-order organization; it therefore supports an edge-organization interpretation without isolating homophily as a causal mechanism. The two experiments together expose the central gap: useful graph signal can be substantial even when simple statistics fail to characterize when a graph model should be used.

The paper makes three contributions:
\begin{enumerate}
    \item \textbf{A frozen graph-vs-MLP decision benchmark.} We evaluate transparent diagnostics---with label-dependent graph statistics restricted to the training split and no test-label access---as explicit actions under fixed asymmetric model portfolios, matched per-architecture tuning grids, common splits and seeds, a practical margin, and test-once accounting.
    \item \textbf{A decision-oriented negative result.} Across 110 units, simple diagnostics do not improve on trivial policies in a stable way; regret and coverage reveal failures that correlation or raw accuracy alone can obscure.
    \item \textbf{A mechanism--decision separation.} A degree-preserving paired intervention shows that diagnostic failure cannot be reduced to graph irrelevance, while a scoped discriminability calculation in Section~\ref{sec:snr} provides limited mechanistic context without serving as a selection theorem.
\end{enumerate}

The remainder of the paper positions the study against graph characterization and model selection, defines the frozen decision protocol, reports the prospective benchmark and intervention, and then presents scoped mechanistic context and limitations.
